%%% FIELD stability-statement
Fix an arm and profile.  If the obstacle-set construction is evaluated at two nonnegative continuous capacities \(u_{a,z}\) and \(u'_{a,z}\), then
\[
d_H^{\|\cdot\|_\infty}\!\left(\mathcal S_{a,z}(u),\mathcal S_{a,z}(u')\right)
\leq \|u_{a,z}-u'_{a,z}\|_\infty.
\]
Consequently, for two base-capacity pairs \((b,u)\) and \((b',u')\), with
\[
\rho_b=\max_{a,z}\|b_{a,z}-b'_{a,z}\|_\infty,
\qquad
\rho_u=\max_{a,z}\|u_{a,z}-u'_{a,z}\|_\infty,
\]
the corresponding sharp-curve constructions obey
\[
d_H^{\|\cdot\|_\infty}\!\left(\mathcal T(b,u),\mathcal T(b',u')\right)
\leq 2\rho_b+2\rho_u.
\]
Both inequalities hold without any restriction on the zero-capacity contact sets.
%%% FIELD stability-proof
Put \(\rho=\|u_{a,z}-u'_{a,z}\|_\infty\), and fix \(f\in\mathcal S_{a,z}(u)\).  Define
\[
(C_\rho f)(r):=[f(r)-\rho]_+.
\tag{1}
\]
The map \(x\mapsto[x-\rho]_+\) is nondecreasing and one-Lipschitz.  Since \(f\) is nonnegative, nondecreasing, and \(L\)-Lipschitz by the obstacle-set definition, (1) is also nonnegative and nondecreasing and satisfies
\[
|(C_\rho f)(r)-(C_\rho f)(r')|
\leq |f(r)-f(r')|
\leq L|r-r'|.
\tag{2}
\]
Moreover, \(f\leq u_{a,z}\leq u'_{a,z}+\rho\) and \(u'_{a,z}\geq0\), whence
\[
0\leq C_\rho f\leq u'_{a,z}.
\tag{3}
\]
Thus \(C_\rho f\in\mathcal S_{a,z}(u')\).  Pointwise,
\[
0\leq f(r)-[f(r)-\rho]_+=\min\{f(r),\rho\}\leq\rho,
\tag{4}
\]
so every element of \(\mathcal S_{a,z}(u)\) is within \(\rho\) of \(\mathcal S_{a,z}(u')\).  Interchanging \(u_{a,z}\) and \(u'_{a,z}\) proves the first Hausdorff inequality.  Notice that (1)--(4) remain valid when either capacity vanishes on an interval or at an isolated point.

For the curve claim, fix
\[
g_z=b_{1,z}-b_{0,z}+f_{1,z}-f_{0,z}\in\mathcal T(b,u).
\]
Apply (1) armwise with \(\rho_u\), writing \(f'_{a,z}=C_{\rho_u}f_{a,z}\).  Then \(f'_{a,z}\in\mathcal S_{a,z}(u')\), and
\[
g'_z=b'_{1,z}-b'_{0,z}+f'_{1,z}-f'_{0,z}\in\mathcal T(b',u').
\]
The triangle inequality and (4) give
\[
\|g-g'\|_\infty
\leq \max_z\bigl(\|b_{1,z}-b'_{1,z}\|_\infty
+\|b_{0,z}-b'_{0,z}\|_\infty
+\|f_{1,z}-f'_{1,z}\|_\infty
+\|f_{0,z}-f'_{0,z}\|_\infty\bigr)
\leq2\rho_b+2\rho_u.
\tag{5}
\]
The same construction with the primed and unprimed pairs exchanged gives the reverse directed bound, proving the second Hausdorff inequality.
%%% FIELD inference-construction
Put
\[
\widehat\gamma_z=\frac1n\sum_{i=1}^n\mathbf 1\{Z_i=z\},
\qquad
\varepsilon_{Z,n}=2\sqrt{\frac{|\mathcal Z|}{n\alpha}},
\qquad
e_{h_n}=2(Lh_n+Mh_n^\beta).
\]
Let \(\widehat{\mathcal E}_n\) be the set of triples \((b',u',p')\in\widehat{\mathcal B}_{\eta,n}\) such that, for every arm and profile, \(u'_{a,z}\geq0\), \(p'_z\geq0\), \(\int_0^1p'_z(r)\,dr=1\), \(c_p\leq p'_z\leq C_p\), and
\[
\max_{a,z}\{\|b'_{a,z}\|_{\mathcal H^\beta},\|u'_{a,z}\|_{\mathcal H^\beta}\}
\vee\max_z\|p'_z\|_{\mathcal H^\beta}\leq M.
\]
Define the finite-multinomial region
\[
\widehat\Gamma_n=
\left\{\gamma\in[0,1]^{|\mathcal Z|}:\sum_z\gamma_z=1,
\ \sum_z|\gamma_z-\widehat\gamma_z|\leq\varepsilon_{Z,n}\right\}
\]
and the preliminary baseline-law region
\[
\widehat{\mathcal P}^{\,0}_{X,n}
=\left\{Q:Q(\{z\}\times dr)=\gamma_zp'_z(r)\,dr
\text{ for some }(b',u',p')\in\widehat{\mathcal E}_n,
\ \gamma\in\widehat\Gamma_n\right\}.
\]
Set \(\widehat{\mathcal P}_{X,n}=\widehat{\mathcal P}^{\,0}_{X,n}\) when this set is nonempty, and otherwise use the observable fallback singleton containing \(Q^0(\{z\}\times dr)=|\mathcal Z|^{-1}dr\).

Let
\[
\mathcal K_{\mathsf R,M,L,\beta}
=\left\{g\in\mathbb C_{\mathcal Z}:\|g\|_\infty\leq\mathsf R,
\ |g_z(r)-g_z(r')|
\leq2M|r-r'|^\beta+2L|r-r'|\ \forall z,r,r'\right\}.
\]
For \((b',u',p')\in\widehat{\mathcal E}_n\), let \(\mathcal T_{h_n}(b',u')\) denote the gridded-curve construction evaluated at \((b',u')\).  Define
\[
\widehat H_n=
\overline{\operatorname{co}}\!\left(
\bigcup_{(b',u',p')\in\widehat{\mathcal E}_n}
\left\{g+e:g\in\mathcal T_{h_n}(b',u'),\ e\in\mathbb C_{\mathcal Z},\
\|e\|_\infty\leq e_{h_n}\right\}\right)
\cap\mathcal K_{\mathsf R,M,L,\beta}.
\]
The band-robust whole-set confidence set is \(\widehat{\mathcal T}_{n,h_n}^{+}=\widehat H_n\) when \(\widehat H_n\neq\varnothing\), and is the fallback singleton \(\{0\}\) otherwise.  Every operation in this construction uses only the observed sample, the externally supplied nuisance band, and predeclared constants; it does not differentiate or hold fixed an obstacle contact face.
%%% FIELD leaf-interval-construction
For every \(G\in\mathcal G\), let
\[
\widehat{\mathcal A}_{G,n}
=\left\{(g,Q):g\in\widehat{\mathcal T}_{n,h_n}^{+},\
Q\in\widehat{\mathcal P}_{X,n},\ Q(G)\geq p_{\min,n}/2\right\}.
\]
If \(\widehat{\mathcal A}_{G,n}\neq\varnothing\), define
\[
\widehat I_G=
\left[
\inf_{(g,Q)\in\widehat{\mathcal A}_{G,n}}\ell_G(g;Q),
\ \sup_{(g,Q)\in\widehat{\mathcal A}_{G,n}}\ell_G(g;Q)
\right].
\]
If \(\widehat{\mathcal A}_{G,n}=\varnothing\), define the observable fallback \(\widehat I_G=[-\mathsf R,\mathsf R]\).  Intervals are therefore constructed for the entire predeclared class; the theorem supplies simultaneous guarantees for every mass-eligible \(G\in\mathcal G_n\).
%%% FIELD main-statement
Let \(\widehat{\mathcal T}_{n,h_n}^{+}\), \(\widehat{\mathcal P}_{X,n}\), and \(\{\widehat I_G:G\in\mathcal G\}\) be the band-robust constructions above, and define
\[
E_{\mathcal T,n}=8r_n+2(Lh_n+Mh_n^\beta),
\qquad
\varepsilon_{Z,n}=2\sqrt{\frac{|\mathcal Z|}{n\alpha}},
\qquad
\Delta_{X,n}=2(r_n+\varepsilon_{Z,n}).
\]
Uniformly over \(P\in\mathfrak P_n^{\mathrm{band}}(r_n)\), with probability at least \(1-\alpha/2-o(1)\), all of the following statements hold simultaneously:
\[
\mathcal T(P_{\mathrm{obs}})\subseteq\widehat{\mathcal T}_{n,h_n}^{+},
\qquad
P_X\in\widehat{\mathcal P}_{X,n},
\]
\[
\sup_{\widehat g\in\widehat{\mathcal T}_{n,h_n}^{+}}
\inf_{g\in\mathcal T(P_{\mathrm{obs}})}\|\widehat g-g\|_\infty
\leq E_{\mathcal T,n},
\qquad
\sup_{Q\in\widehat{\mathcal P}_{X,n}}\|Q-P_X\|_{\mathrm{TV}}
\leq\Delta_{X,n},
\]
and, writing \(\widehat I_G=[\widehat{\underline\theta}_G,\widehat{\overline\theta}_G]\), for every \(G\in\mathcal G_n\),
\[
[\underline\theta_G,\overline\theta_G]\subseteq\widehat I_G,
\qquad
0\leq
\max\{\underline\theta_G-\widehat{\underline\theta}_G,
\widehat{\overline\theta}_G-\overline\theta_G\}
\leq E_{\mathcal T,n}+\frac{4\mathsf R\Delta_{X,n}}{p_{\min,n}}.
\]
Thus the whole-set directed excess is \(O(r_n)\) on this event under the declared analysis mesh and tends to zero because \(r_n\to0\).  The conservative simultaneous leaf excess tends to zero along any sequence satisfying \((r_n+n^{-1/2})/p_{\min,n}\to0\).  For every data-measurable \(\widehat G\in\mathcal G_n\), the same event implies \([\underline\theta_{\widehat G},\overline\theta_{\widehat G}]\subseteq\widehat I_{\widehat G}\), without conditioning on its selection rule.

The score-envelope, score-moment, and score-entropy assumptions alone do not imply validity of a derivative-aware common multiplier critical value or the sharper local-mass leaf rate \((np_{\min,n})^{-1/2}\): they admit a zero score class whose multiplier process is identically zero while an eligible empirical leaf mass has nonzero sampling error.  Obtaining that sharper result requires an additional, explicitly declared uniform score-to-statistic expansion and a conditional multiplier approximation.  Hence changing contacts do not obstruct the proved band-robust whole-set result, whereas efficient local-mass multiplier calibration is not identified by the present premises.
%%% FIELD main-proof
Write \(\eta^\circ=(b^\circ,u^\circ,p^\circ)\) for the true observed-law nuisance, and let
\[
\mathcal E_{B,n}=
\left\{\eta^\circ\in\widehat{\mathcal B}_{\eta,n},\
\operatorname{diam}_\infty(\widehat{\mathcal B}_{\eta,n})\leq2r_n\right\}.
\tag{1}
\]
The external-band prerequisite and the definition of the band-valid subclass give
\[
\inf_{P\in\mathfrak P_n^{\mathrm{band}}(r_n)}P(\mathcal E_{B,n})
\geq1-\alpha/4-o(1).
\tag{2}
\]
On (1), the score-density and observed-law H\"older assumptions put \(\eta^\circ\) in the admissible portion \(\widehat{\mathcal E}_n\) of the band.

We first prove whole-set coverage.  Put \(e_{h_n}=2(Lh_n+Mh_n^\beta)\) and \(\mathbb B_\infty=\{e\in\mathbb C_{\mathcal Z}:\|e\|_\infty\leq1\}\).  The uniform grid-Hausdorff theorem, evaluated at the true nuisance, gives
\[
\mathcal T(P_{\mathrm{obs}})
\subseteq\mathcal T_{h_n}(b^\circ,u^\circ)+e_{h_n}\mathbb B_\infty.
\tag{3}
\]
Every \(g\in\mathcal T(P_{\mathrm{obs}})\) satisfies \(\|g\|_\infty\leq t_{\max}=\mathsf R\): by law sharpness it is the CATE curve of a compatible law, and bounded event times put both conditional arm means in \([0,t_{\max}]\).  Also, the observed-law H\"older restriction and the residual Lipschitz restriction imply
\[
|g_z(r)-g_z(r')|
\leq2M|r-r'|^\beta+2L|r-r'|.
\tag{4}
\]
Thus \(\mathcal T(P_{\mathrm{obs}})\subseteq\mathcal K_{\mathsf R,M,L,\beta}\).  On \(\mathcal E_{B,n}\), (3), (4), and \(\eta^\circ\in\widehat{\mathcal E}_n\) show
\[
\mathcal T(P_{\mathrm{obs}})\subseteq\widehat H_n
=\widehat{\mathcal T}_{n,h_n}^{+}.
\tag{5}
\]
In particular the fallback clause is inactive.

For clarity, the set in (5) is compact even though the displayed expansion before clipping uses a sup-norm ball.  Indeed, \(\mathcal K_{\mathsf R,M,L,\beta}\) is closed, uniformly bounded, and has a common modulus of continuity.  To verify compactness without invoking an undeclared theorem, take any sequence in this class, extract a diagonal subsequence converging on a countable dense subset of \([0,1]\) in each of the finitely many profiles, and use the common modulus to make that subsequence uniformly Cauchy on a finite score net.  Its uniform limit is continuous and retains the defining inequalities.  Hence \(\mathcal K_{\mathsf R,M,L,\beta}\) is sequentially compact in the sup norm and therefore compact.  The common-modulus inequalities also show that this class is convex.  The set \(\widehat H_n\) is the intersection of it with a closed convex set, so it is compact and convex; the fallback singleton has the same properties.

We next obtain a contact-set-free excess bound.  Fix \((b',u',p')\in\widehat{\mathcal E}_n\).  On \(\mathcal E_{B,n}\), band diameter gives
\[
\max_{a,z}\|b'_{a,z}-b^\circ_{a,z}\|_\infty\leq2r_n,
\qquad
\max_{a,z}\|u'_{a,z}-u^\circ_{a,z}\|_\infty\leq2r_n.
\tag{6}
\]
The obstacle-stability theorem therefore yields
\[
d_H\{\mathcal T(b',u'),\mathcal T(b^\circ,u^\circ)\}\leq8r_n.
\tag{7}
\]
The gridded residual interpolants built from \(u'\) lie in the continuous obstacle for \(u'\).  In fact, on a cell \([r_j,r_{j+1}]\), each endpoint satisfies
\[
x_k\leq[u'_{a,z}(r_k)-Mh_n^\beta]_+\leq u'_{a,z}(r)
\quad(k=j,j+1),
\tag{8}
\]
because \(|u'_{a,z}(r_k)-u'_{a,z}(r)|\leq Mh_n^\beta\) and feasibility of a nonnegative \(x_k\) forces the positive-part upper bound to be nonnegative.  Linear interpolation preserves (8); its slopes lie in \([0,L]\).  Thus
\[
\mathcal T_{h_n}(b',u')\subseteq\mathcal T(b',u').
\tag{9}
\]
Equations (7)--(9) imply that every element of every expanded set inside the union defining \(\widehat H_n\) is within \(8r_n+e_{h_n}\) of \(\mathcal T(P_{\mathrm{obs}})\).  The true set is closed and convex by the law-sharpness theorem.  Distance to a convex set is convex: for finitely many \(y_k\), choose \(t_k\) in the true set arbitrarily close to each distance and observe that \(\sum_k\lambda_kt_k\) remains in the true set.  Distance is also one-Lipschitz, so passage to the closed convex hull preserves the same bound.  Intersecting with \(\mathcal K_{\mathsf R,M,L,\beta}\) cannot increase the supremum.  Therefore, on \(\mathcal E_{B,n}\),
\[
\sup_{\widehat g\in\widehat{\mathcal T}_{n,h_n}^{+}}
\inf_{g\in\mathcal T(P_{\mathrm{obs}})}\|\widehat g-g\|_\infty
\leq8r_n+2(Lh_n+Mh_n^\beta)=E_{\mathcal T,n}.
\tag{10}
\]
No contact-set, active-face, or differentiability condition entered (6)--(10).

We now construct baseline-law coverage.  Let \(\gamma_{0,z}=P_X(Z=z)\) and retain \(\widehat\gamma_z\) from the definition.  Independent sampling gives
\[
E_P\|\widehat\gamma-\gamma_0\|_2^2
=\sum_z\operatorname{Var}_P(\widehat\gamma_z)
=\frac{1-\sum_z\gamma_{0,z}^2}{n}
\leq\frac1n.
\tag{11}
\]
Since \(\|x\|_1\leq\sqrt{|\mathcal Z|}\|x\|_2\), and since \(E[Y]\geq tP(Y>t)\) for every nonnegative \(Y\), equation (11) gives
\[
P\!\left\{\|\widehat\gamma-\gamma_0\|_1>
2\sqrt{\frac{|\mathcal Z|}{n\alpha}}\right\}
\leq
P\!\left\{\|\widehat\gamma-\gamma_0\|_2^2>\frac4{n\alpha}\right\}
\leq\frac\alpha4.
\tag{12}
\]
Call the complementary event \(\mathcal E_{Z,n}\).  On \(\mathcal E_{B,n}\cap\mathcal E_{Z,n}\), the true pair \((\gamma_0,p_0)\) is admitted by the preliminary baseline-law region, hence
\[
P_X\in\widehat{\mathcal P}_{X,n}.
\tag{13}
\]

Define the total-variation norm by
\[
\|Q-P_X\|_{\mathrm{TV}}
=\sup_{|\varphi|\leq1}\left|\int\varphi\,d(Q-P_X)\right|.
\tag{14}
\]
Every \(Q\in\widehat{\mathcal P}_{X,n}\) on the joint event has the form \(Q(\{z\}\times dr)=\gamma_zp'_z(r)\,dr\).  The triangle inequality inside the two probability simplices, band diameter, and the normalization of each conditional density give
\[
\begin{aligned}
\|Q-P_X\|_{\mathrm{TV}}
&\leq\sum_z|\gamma_z-\gamma_{0,z}|\int_0^1p'_z(r)\,dr
+\sum_z\gamma_{0,z}\int_0^1|p'_z(r)-p^\circ_z(r)|\,dr\\
&\leq2\varepsilon_{Z,n}+2r_n
=\Delta_{X,n}.
\end{aligned}
\tag{15}
\]
Here the factor two in the first term follows because both \(\gamma\) and \(\gamma_0\) lie in the \(\varepsilon_{Z,n}\)-ball about \(\widehat\gamma\).

Combining (2) and (12) by the union bound proves that (5), (10), (13), and (15) hold simultaneously with probability at least \(1-\alpha/2-o(1)\), uniformly over the certified subclass.

It remains to verify the leaf assertions.  Fix \(G\in\mathcal G_n\).  By eligibility, \(P_X(G)\geq p_{\min,n}>0\).  Equations (5) and (13) show that every \(g\in\mathcal T(P_{\mathrm{obs}})\), paired with \(P_X\), is feasible for the interval program because \(P_X(G)\geq p_{\min,n}/2\).  Taking the infimum and supremum therefore gives
\[
[\underline\theta_G,\overline\theta_G]\subseteq\widehat I_G.
\tag{16}
\]

For the excess bound, take any feasible \((\widehat g,Q)\).  By (10), for every \(\epsilon>0\) there is \(g^0\in\mathcal T(P_{\mathrm{obs}})\) with \(\|\widehat g-g^0\|_\infty\leq E_{\mathcal T,n}+\epsilon\).  Put \(D=Q(G)\), \(D_0=P_X(G)\), and \(N_Q=\int_Gg^0\,dQ\), \(N_0=\int_Gg^0\,dP_X\).  Since \(D\geq p_{\min,n}/2\), \(D_0\geq p_{\min,n}\), and \(\|g^0\|_\infty\leq\mathsf R\), (14)--(15) imply
\[
\begin{aligned}
|\ell_G(\widehat g;Q)-\ell_G(g^0;P_X)|
&\leq\|\widehat g-g^0\|_\infty
+\frac{|N_Q-N_0|}{D}
+\frac{|N_0|\,|D-D_0|}{DD_0}\\
&\leq E_{\mathcal T,n}+\epsilon
+\frac{2\mathsf R\Delta_{X,n}}{D}\\
&\leq E_{\mathcal T,n}+\epsilon
+\frac{4\mathsf R\Delta_{X,n}}{p_{\min,n}}.
\end{aligned}
\tag{17}
\]
Letting \(\epsilon\downarrow0\), then taking the infimum and supremum over feasible pairs in (17), proves the simultaneous endpoint-excess inequality.  Because (16) holds for every eligible leaf on one event, substituting any data-measurable \(\widehat G\in\mathcal G_n\) proves the post-selection assertion without conditioning.

The analysis-mesh definition gives \(Lh_n+Mh_n^\beta=O(r_n)\).  Hence (10) is \(O(r_n)\) and tends to zero under \(r_n\to0\).  Since \(|\mathcal Z|\) and \(\alpha\) are fixed, \(\varepsilon_{Z,n}=O(n^{-1/2})\); (17) therefore shrinks under the stated condition \((r_n+n^{-1/2})/p_{\min,n}\to0\).  This is deliberately conservative relative to the desired local-mass term.

Finally, we prove that the remaining derivative-aware claim is not a consequence of the current score premises.  Consider the following allowed instance, taking admissible constants with \(M\geq1\) and \(c_p\leq1\leq C_p\).  Let \(\mathcal Z=\{z_0\}\), let \(R\) be uniform on \([0,1]\), draw \(A\) independently with probability \(1/2\) in each arm, and set \(T(0)=T(1)=0\) and \(C(0)=C(1)=t_{\max}+1\).  Then \(p_{z_0}=1\), \(b_{a,z_0}=u_{a,z_0}=0\), treatment overlap holds because \(\varepsilon<1/2\), uncensored overlap holds with probability one, and all residual and observed-law smoothness conditions hold.  Use the predeclared leaf \(G=\{(z_0,r):r\leq1/2\}\), with \(p_{\min,n}\leq1/2\); its singleton class has finite VC dimension.  An oracle singleton nuisance band satisfies the external band prerequisites.  For every index in \(\mathcal D_n\times\mathcal G_n\), set \(\psi\equiv0\) and take \(F_\psi\equiv0\).  The class is centered, its \((2+\delta_\psi)\)-moment is zero, and its covering number is one, so the score-envelope, score-moment, and score-entropy assumptions all hold.  Nevertheless, for every odd \(n\),
\[
\widehat P_n(G)=\frac1n\sum_{i=1}^n\mathbf1\{R_i\leq1/2\}\neq\frac12=P_X(G)
\quad\text{for every sample},
\tag{18}
\]
whereas every multiplier process formed from the declared zero scores, and hence every one-sided or two-sided multiplier quantile of that process, is identically zero.  Thus those score premises alone cannot calibrate even the nonzero leaf-mass error in (18), much less an optimized endpoint through a changing critical cone.  An explicit uniform score-to-statistic expansion with a controlled remainder and a conditional multiplier approximation are logically additional premises.  This counterexample establishes the claimed obstruction while (1)--(17) close the contact-stable whole-set and conservative leaf parts.
